\documentclass[UTF8,fontset=none,11pt,a4paper]{ctexart}
\usepackage{ipara-handbook}
\usepackage{amsmath,booktabs,tabularx,hyperref}
\handbooksetup{OS-B03 Process × File Reference}{408 · OS Internal Bridge}{FD · OFD · Inode · Lifetime}
\begin{document}
\handbooktitle{Process × File Reference}{文件描述符不是文件：引用层次怎样决定 fork、dup、close 和 unlink}{408 · OS-B03}
\begin{corebox}{Mother Interface}
\[
\boxed{\text{Task}\to\text{FD Table Entry}\to\text{Open-File Description}\to\text{Inode/File Object}\to\text{Reference/Lifetime Change}}
\]
本 Bridge 拥有引用关系在不同层级的交接，不重新定义路径解析或文件内容读写。
\end{corebox}
\begin{methodbox}{Owners 与边界}
OS01 Owns task 与 per-process fd table；OS05 Owns pathname、inode、open-file state 与生命周期。本册 Owns“一个操作改变哪一层引用”；I/O 与 Page Cache 进入 OS-B04/X-I02。
\end{methodbox}
\section{三层对象必须分开}
fd 是进程表中的整数索引；fd table entry 指向 open-file description，后者保存当前 offset、status flags 等打开状态；多个 description 可指向同一 inode/file object。fork/dup 复制的是指向关系，不等同于复制文件内容。
\section{操作对照}
\begin{center}\small\begin{tabularx}{\linewidth}{lXX}\toprule 操作&主要新增/共享关系&常见后果\\\midrule fork&子表复制，条目指向同一 OFD&父子共享 offset\\ dup&同一进程新增 fd，指向同一 OFD&close 一个不必关闭 OFD\\ unlink&移除目录名，保留已有引用&最后引用消失后才可回收\\ close&移除一个 fd table entry&可能触发 OFD/inode 引用递减\\\bottomrule\end{tabularx}\end{center}
\section{最小例子与 Anti-Bridge}
父进程打开文件后 fork，父子交替 read；由于共享 OFD，文件 offset 可能共同推进。若重新 open，则得到新的 OFD，offset 不共享。\code{fd} $\neq$ inode；unlink $\neq$ 立即删除已有打开对象；close $\neq$ 必然销毁文件内容。
\begin{corebox}{Compression}
题目给的是哪个对象？操作复制的是索引、OFD 还是 inode 引用？哪个字段共享？最后一个引用何时消失？
\end{corebox}
\end{document}
